% Answer Key: Prompt Engineering — Quiz (20 MCQs)
\documentclass[11pt,a4paper]{article}
\usepackage[utf8]{inputenc}
\usepackage{geometry}
\geometry{margin=1in}
\begin{document}
\begin{center}
  {\LARGE Answer Key — Prompt Engineering Quiz}\\[6pt]
\end{center}
\vspace{0.5cm}
\begin{enumerate}
  \item B — The textual instruction, question, or example you give an AI system.
  \item B — They predict the next token based on patterns learned from large text corpora.
  \item A — True; tokenizers commonly split long words into sub-word tokens.
  \item A — Self-attention compares and weights other tokens so representations become context-aware.
  \item C — Embeddings are numeric vectors in a meaning space.
  \item B — Hallucination: plausible-sounding but incorrect or fabricated information.
  \item B — False; larger size helps but does not guarantee no hallucinations or errors.
  \item D — All of the above (role, format, audience are useful levers).
  \item A — Query/Key/Value mapping: Query = what I'm looking for; Key = what I contain; Value = what I provide.
  \item B — Fine-tuning: additional training on curated examples to specialize behavior.
  \item B — Emergent ability example: few-shot learning appearing in larger models.
  \item A — Position embeddings encode token order since attention is permutation-invariant.
  \item C — Unlimited context-window is NOT a typical capability; it's a limitation when claiming it's unlimited.
  \item A — Tokenization allows handling rare or long words by splitting into subparts.
  \item B — Example of effective constrained prompt is the teacher + format + length example.
  \item A — Embedding metric: high-dimensional numeric vector; nearby vectors = semantic similarity.
  \item A — Instruction tuning teaches the model to follow human instructions better.
  \item A — Verifying facts with reliable external sources or retrieval helps reduce hallucinations.
  \item A — Multi-head attention allows attending to multiple representation subspaces and capture different relationships.
  \item A — Purpose, context, and constraints are the three elements good prompts typically provide.
\end{enumerate}

\vspace{0.5cm}
  extbf{Notes for graders:} Each MCQ = 1 point. For any ambiguous phrasing, use the explanation above as the authoritative key; award partial credit only if a student shows clear correct reasoning in comments.
\end{document}
